\documentclass[../../main]{subfiles}

\begin{document}

\section{Reticulados}
\label{sec:matematica:relacoes:reticulados}

\begin{defn}[Majorante]
  \label{def:matematica:relacoes:reticulados:majorante}

  Seja ((P,\le)) um conjunto parcialmente ordenado e seja
  (A\subseteq P).

  Um elemento (u\in P) é denominado majorante de (A) quando

  [
    (\forall a\in A)
    (a\le u).
  ]

\end{defn}

\begin{defn}[Minorante]
  \label{def:matematica:relacoes:reticulados:minorante}

  Seja (A\subseteq P).

  Um elemento (l\in P) é denominado minorante de (A) quando

  [
    (\forall a\in A)
    (l\le a).
  ]

\end{defn}

\begin{defn}[Supremo]
  \label{def:matematica:relacoes:reticulados:supremo}

  Seja (A\subseteq P).

  O supremo de (A), quando existe, é o menor dos majorantes de (A).

  É denotado por

  [
    \sup(A).
  ]

\end{defn}

\begin{defn}[Ínfimo]
  \label{def:matematica:relacoes:reticulados:infimo}

  Seja (A\subseteq P).

  O ínfimo de (A), quando existe, é o maior dos minorantes de (A).

  É denotado por

  [
    \inf(A).
  ]

\end{defn}

\begin{defn}[Junção]
  \label{def:matematica:relacoes:reticulados:juncao}

  Sejam (a,b\in P).

  A junção de (a) e (b), quando existe, é definida por

  [
    a\vee b
    =
    \sup{a,b}.
  ]

\end{defn}

\begin{defn}[Encontro]
  \label{def:matematica:relacoes:reticulados:encontro}

  Sejam (a,b\in P).

  O encontro de (a) e (b), quando existe, é definido por

  [
    a\wedge b
    =
    \inf{a,b}.
  ]

\end{defn}

\begin{defn}[Reticulado]
  \label{def:matematica:relacoes:reticulados:reticulado}

  Um reticulado é um conjunto parcialmente ordenado
  ((P,\le)) tal que, para quaisquer elementos
  (a,b\in P),

  [
    a\vee b
  ]

  e

  [
    a\wedge b
  ]

  existem.

\end{defn}

\begin{prop}[Idempotência]

  Para quaisquer (a,b\in P),

  [
    a\vee a=a
  ]

  e

  [
    a\wedge a=a.
  ]

\end{prop}

\begin{proof}

  Decorre diretamente da definição de supremo e ínfimo.

\end{proof}

\begin{prop}[Comutatividade]

  Para quaisquer (a,b\in P),

  [
    a\vee b
    =
    b\vee a
  ]

  e

  [
    a\wedge b
    =
    b\wedge a.
  ]

\end{prop}

\begin{proof}

  Os conjuntos ({a,b}) e ({b,a}) coincidem.

\end{proof}

\begin{prop}[Associatividade]

  Para quaisquer (a,b,c\in P),

  [
    (a\vee b)\vee c
    =
    a\vee(b\vee c)
  ]

  e

  [
    (a\wedge b)\wedge c
    =
    a\wedge(b\wedge c).
  ]

\end{prop}

\begin{prop}[Leis de Absorção]

  Para quaisquer (a,b\in P),

  [
    a\vee(a\wedge b)
    =
    a
  ]

  e

  [
    a\wedge(a\vee b)
    =
    a.
  ]

\end{prop}

\begin{proof}

  Segue das propriedades universais de supremo e ínfimo.

\end{proof}

\begin{ex}

  O conjunto das partes

  [
    \mathcal P(A)
  ]

  munido da inclusão

  [
    \subseteq
  ]

  forma um reticulado.

  Nesse caso,

  [
    X\vee Y
    =
    X\cup Y
  ]

  e

  [
    X\wedge Y
    =
    X\cap Y.
  ]

\end{ex}

\end{document}
